\newpage{\ }
\thispagestyle{empty}
\chapter*{Resumen}

\noindent Para muchas personas con dificultades graves del habla, comunicar una idea tan simple como pedir agua pasa por seleccionar una secuencia de pictogramas en una aplicación. Es una situación frecuente entre quienes conviven con un Trastorno del Espectro Autista (TEA), entre otras condiciones. Las aplicaciones disponibles en español suelen limitarse a reproducir esa secuencia tal cual, lo que produce construcciones artificiales del tipo ``yo querer agua''.

Este trabajo presenta Dixi, un sistema que toma como entrada secuencias de pictogramas del catálogo ARASAAC (Centro Aragonés para la Comunicación Aumentativa y Alternativa) y produce frases naturales en español. La meta es que el mensaje resultante se aproxime a cómo lo formularía un hablante sin barreras comunicativas, de modo que la conversación deje de depender del esfuerzo del interlocutor por reconstruir lo que el usuario quiere decir.

Para alcanzar ese objetivo se exploran tres caminos distintos. El primero parte de reglas gramaticales explícitas que describen cómo se construye una frase correcta en español, escritas con la herramienta \textit{Lark}. El segundo se apoya en la semántica que el lenguaje deja en grandes colecciones de texto, capturada mediante \textit{embeddings} de \textit{FastText} y un \textit{pipeline} propio que selecciona y combina las palabras adecuadas. El tercero recurre a un modelo de lenguaje de gran escala (\textit{LLM}), \textit{GPT-4o}, guiado mediante ingeniería de \textit{prompts}. Sobre estas tres aproximaciones se construyen además tres motores híbridos, que combinan las virtudes de unas y otras para corregir sus puntos débiles.

Comparar técnicas tan distintas exige una evaluación rigurosa, que aquí se aborda por dos vías complementarias. La primera consiste en un conjunto de métricas automáticas calculadas sobre un \textit{golden standard} de 164 casos construido a partir de una encuesta a hablantes nativos, con varias frases de referencia por caso. La segunda es una encuesta de elección forzada en la que personas reales eligen, sin saber qué motor produjo cada frase, cuál les suena más natural. El \textit{LLM} encabeza ambas evaluaciones: concentra el 44\% de las preferencias humanas y mantiene una ventaja estadísticamente significativa frente a la gramática y a los \textit{embeddings}.

Cerrado el ciclo de evaluación, el sistema se entrega al usuario final mediante una aplicación móvil multiplataforma desarrollada con React Native y Expo. La aplicación organiza los pictogramas por categorías semánticas, permite construir frases mediante selección directa y convierte el resultado en habla a través de un motor de síntesis de voz, completando así el recorrido desde la idea hasta el mensaje hablado.

\vspace{.5cm}

\noindent \textit{Palabras clave:} generación de lenguaje natural, sistemas aumentativos y alternativos de comunicación (SAAC), pictogramas ARASAAC, Trastorno del Espectro Autista (TEA), gramáticas libres de contexto, \textit{embeddings}, modelos de lenguaje de gran escala, hibridación, síntesis de voz, evaluación humana.
